\documentclass[12pt]{article}
\usepackage[margin=1in]{geometry}
\usepackage{graphicx}
\usepackage{amsmath}
\usepackage{natbib}
\usepackage{booktabs}
\usepackage{hyperref}
\usepackage{xcolor}
\usepackage{textcomp}
\usepackage{longtable}
\usepackage{array}
\usepackage[T1]{fontenc}

\bibliographystyle{unsrtnat}

\title{A SUMO E3 Ligase Promotes Long Non-Coding RNA Transcription to Regulate Small RNA-Directed DNA Elimination}
\author{}
\date{}

\begin{document}
\maketitle

%% ============================================================
%% ABSTRACT
%% ============================================================
\begin{abstract}
Small RNAs target complementary chromatin regions for gene silencing through nascent long non-coding RNAs (lncRNAs).
In the ciliated protozoan \textit{Tetrahymena thermophila}, Piwi-associated small RNAs (scnRNAs) interact with nascent lncRNA transcripts from the somatic macronuclear (MAC) genome, inducing target-directed small RNA degradation (TDSD).
The scnRNAs that escape TDSD subsequently direct programmed DNA elimination of germline-limited internal eliminated sequences (IESs) in the developing new MAC.
Here we demonstrate that the SUMO E3 ligase Ema2 is required for the accumulation of lncRNAs from the parental MAC genome and consequently for TDSD and completing DNA elimination to produce viable sexual progeny.
Ema2 directly interacts with the SUMO E2 conjugating enzyme Ubc9 and enhances SUMOylation of the conserved transcription elongation factor Spt6.
Cell fractionation experiments reveal that Ema2 promotes the chromatin association of Spt6 and RNA polymerase II (RNAPII) in the parental MAC.
However, a SUMOylation-defective Spt6 mutant (Spt6-C-KR) retains chromatin association and supports lncRNA transcription, although it exhibits a mild DNA elimination defect.
These findings indicate that Ema2-directed SUMOylation actively promotes lncRNA transcription, providing a prerequisite for genome--small RNA communication during programmed DNA elimination.
\end{abstract}

%% ============================================================
%% 1. INTRODUCTION
%% ============================================================
\section{Introduction}

Small RNAs of approximately 20--30 nucleotides that are complexed with Argonaute family proteins target either mRNAs for post-transcriptional silencing or chromatin regions for transcriptional gene silencing~\citep{Holoch2015, Wilson2013}.
For the latter process, small RNAs are generally considered to recognize their genomic targets via nascent long non-coding RNAs (lncRNAs) to induce heterochromatin formation.
Therefore, chromatin regions targeted for silencing by small RNAs must paradoxically be transcribed to provide nascent lncRNAs.

In fission yeast, small interfering RNAs (siRNAs) mediate the deposition of histone H3 lysine 9 di- and tri-methylation (H3K9me2/3) for heterochromatin assembly at centromeric repeats~\citep{Volpe2002, Hall2002}.
During the M-to-S phase transition, phosphorylation of histone H3 serine 10 evicts the HP1 protein Swi6, passively permitting lncRNA transcription and subsequent H3K9me2/3 deposition on newly assembled nucleosomes~\citep{Chen2008, Kloc2008}.
In contrast, in \textit{Drosophila}, the HP1 paralog Rhino actively recruits dedicated transcription factors to piRNA clusters, enabling lncRNA transcription from heterochromatin while preventing mRNA transcription from embedded transposons~\citep{LeThomas2014, Andersen2017}.
A similar active heterochromatin-dependent lncRNA transcription occurs in plants, where SHH1 reads H3K9 methylation marks and recruits the plant-specific RNA polymerase IV~\citep{Law2013, Law2011}.

Because small RNA-producing loci are also small RNA targets in most studied systems, it is challenging to separately investigate lncRNA transcription for small RNA biogenesis and that for downstream effector recruitment.
The ciliated protozoan \textit{Tetrahymena thermophila} offers a unique experimental system in which the source and target loci of small RNAs reside in different nuclei during programmed DNA elimination~\citep{Chalker2013, Noto2017}.

Each \textit{Tetrahymena} cell harbors a diploid germline micronucleus (MIC) and a polyploid somatic macronucleus (MAC).
During conjugation, the MIC undergoes meiosis and fertilization, followed by the formation of the new MIC and MAC, while the parental MAC is degraded~\citep{Chalker2013}.
Programmed DNA elimination in the new MAC downsizes the ${\sim}200$~Mb MIC genome to the ${\sim}103$~Mb MAC genome by removing ${\sim}12{,}000$ internal eliminated sequences (IESs), many of which include transposons~\citep{Noto2017}.

This process is regulated by three types of lncRNAs produced in different nuclei at distinct stages~\citep{Aronica2008, Schoeberl2011}.
During meiotic prophase, bidirectional lncRNAs transcribed genome-wide in the MIC by RNA polymerase II are processed into ${\sim}29$-nt scnRNAs by Dicer homologs and loaded onto the Piwi protein Twi1~\citep{Chalker2001, Schoeberl2012, Mochizuki2004b, Malone2005, Mochizuki2005, Mochizuki2002, Noto2010}.
The Piwi--scnRNA complex then moves into the parental MAC, where parental macronuclear lncRNAs (pMAC-lncRNAs) are transcribed bidirectionally during mid-conjugation stages~\citep{Woo2016}.
The RNA helicase Ema1 promotes the base-pairing interaction between pMAC-lncRNAs and scnRNAs~\citep{Aronica2008}, triggering target-directed small RNA degradation (TDSD), a process reminiscent of target-directed microRNA degradation~\citep{Han2023}.
TDSD results in selective retention of IES-derived scnRNAs that subsequently direct DNA elimination in the new MAC~\citep{Aronica2008, Schoeberl2012, Mochizuki2004a, Noto2015, Noto2018}.

Among the three types of lncRNAs, pMAC-lncRNA does not produce small RNAs and is specialized for receptor function.
Therefore, pMAC-lncRNA transcription and the subsequent TDSD in the parental MAC provide a unique paradigm for investigating how the genome is transcribed to communicate with small RNAs.
In the present study, we show that Ema2, a conjugation-specific SUMO E3 ligase, is required for pMAC-lncRNA transcription in \textit{Tetrahymena} and thus provides a molecular tool to dissect the role and mechanism of this lncRNA transcription.

%% ============================================================
%% 2. RESULTS
%% ============================================================
\section{Results}

\subsection{Ema2 Is Exclusively Expressed During Conjugation and Localized in the MAC}

As part of a systematic investigation into genes highly upregulated during conjugation~\citep{Loidl2021}, we explored the function of EMA2 (TTHERM\_00113330).
EMA2 mRNA is exclusively expressed during conjugation.
The encoded Ema2 protein, tagged with HA at the endogenous locus, was undetectable in vegetative cells and first appeared in the MAC at 3 hours post-induction of mating (hpm).
By 6~hpm, Ema2-HA remained localized in the parental MAC.
At the onset of new MAC development (8~hpm), Ema2 disappeared from the parental MAC and appeared in the new MAC, subsequently fading by 12--14~hpm.
This conjugation-specific expression and MAC-to-new-MAC localization switch are reminiscent of factors involved in DNA elimination, such as the Piwi protein Twi1 and PRC2~\citep{Mochizuki2002, Liu2007, Noto2010}.

\subsection{Ema2 Is Required for Completing DNA Elimination}

DNA elimination in exconjugants at 36~hpm was analyzed by DNA fluorescent in situ hybridization (FISH) using probes complementary to the transposable element Tlr1~\citep{Wuitschick2002}.
DNA elimination is normally completed by ${\sim}14$--18~hpm in wild-type cells~\citep{Austerberry1984, Mutazono2019}.
The Tlr1 element was detected only in the MICs in exconjugants from wild-type crosses, whereas in EMA2 somatic knockout (KO) exconjugants, Tlr1 was detected in both MICs and MACs.

The IES retention index was calculated to quantify the extent of DNA elimination failure.
The FISH signal intensity in new MACs was lower in EMA2 KO exconjugants than in TWI1 KO exconjugants, in which DNA elimination is completely blocked~\citep{Noto2015}.
Statistical analysis by the Welch two-sample $t$-test confirmed that the difference in IES retention index between EMA2 KO and wild-type cells was highly significant ($p < 0.001$), as was the difference between EMA2 KO and TWI1 KO ($p < 0.001$).
Thus, DNA elimination was partially blocked in EMA2 KO exconjugants (Table~\ref{tab:phenotype_summary}).

Consistent with the requirement of DNA elimination for progeny viability~\citep{Cheng2010, Vogt2013}, EMA2 KO cells failed to produce viable progeny in two independent crosses (Table~\ref{tab:viability}).
We conclude that maternally expressed EMA2 is required for completing DNA elimination.

\begin{table}[htbp]
\centering
\caption{Progeny viability test. Conjugating pairs were isolated at 6~hpm and tested for viability. ``n'' indicates the total number of pairs tested.}
\label{tab:viability}
\begin{tabular}{lccc}
\toprule
\textbf{Cross} & \textbf{Pairs tested (n)} & \textbf{Viable progeny (\%)} & \textbf{Phenotype} \\
\midrule
Wild-type $\times$ Wild-type & $>$50 & $>$90 & Normal \\
EMA2 KO Cross 1 & $>$50 & 0 & Lethal \\
EMA2 KO Cross 2 & $>$50 & 0 & Lethal \\
\bottomrule
\end{tabular}
\end{table}

\subsection{Ema2 Is Required for Target-Directed Small RNA Degradation (TDSD)}

We next investigated whether EMA2 is involved in TDSD.
Production of scnRNAs occurs from IESs and surrounding MDS regions in the MIC at early conjugation stages (${\sim}2$--3~hpm), and scnRNAs complementary to the MAC genome are subjected to TDSD in the parental MAC during mid-conjugation (${\sim}3$--7~hpm)~\citep{Mochizuki2004a, Aronica2008, Schoeberl2012, Noto2015}.

Northern blot analysis using the 50-nt Mi-9 probe, which is complementary to a repetitive sequence found in both IESs and MDSs~\citep{Aronica2008}, revealed that in wild-type cells, Mi-9-complementary scnRNAs were detected at 3~hpm, reduced at 4.5~hpm, and became undetectable at 6~hpm due to TDSD.
In contrast, Mi-9-complementary scnRNAs remained detectable at 6~hpm in EMA2 KO cells, indicating that Ema2 is required for TDSD.

High-throughput small RNA sequencing at 3 and 8~hpm confirmed the TDSD defect genome-wide.
Sequenced 26--32-nt small RNAs (corresponding to scnRNAs) were mapped to genomic tiles of MDSs (10,235 tiles) and IESs (4,691 Type-A IES tiles)~\citep{Schoeberl2012}.
In wild-type cells, scnRNAs mapping to MDS tiles were greatly reduced from 3 to 8~hpm (Wilcoxon rank sum test, $p < 0.001$), whereas in EMA2 KO cells this reduction was only slight ($p < 0.001$ compared to wild-type at 8~hpm).
This TDSD defect in EMA2 KO cells appeared milder than that in EMA1 KO cells, in which MDS-complementary scnRNA levels remained constant until 8~hpm.
Crucially, the amounts of IES-complementary scnRNAs at 3 and 8~hpm were comparable across all strains ($p > 0.05$), consistent with the established principle that only MDS-complementary scnRNAs are targeted for TDSD (Table~\ref{tab:scnRNA}).

\begin{table}[htbp]
\centering
\caption{Summary of small RNA sequencing results. Normalized scnRNA levels (RPKM) mapped to MDS and IES tiles at 3 and 8~hpm. Statistical significance was assessed by the Wilcoxon rank sum test.}
\label{tab:scnRNA}
\begin{tabular}{lcccc}
\toprule
\textbf{Genotype} & \textbf{MDS 3~hpm} & \textbf{MDS 8~hpm} & \textbf{IES 3~hpm} & \textbf{IES 8~hpm} \\
\midrule
Wild-type & Baseline & Greatly reduced$^{***}$ & Baseline & Comparable (ns) \\
EMA2 KO & Increased & Slightly reduced$^{***}$ & Comparable & Comparable (ns) \\
EMA1 KO & Increased & Unchanged$^{***}$ & Comparable & Comparable (ns) \\
\bottomrule
\multicolumn{5}{l}{\footnotesize $^{***}p < 0.001$ (Wilcoxon rank sum test); ns, $p > 0.05$}
\end{tabular}
\end{table}

\subsection{Ema2 Is Required for Accumulation of SUMOylated Proteins During Conjugation}

Ema2 possesses an SP-RING domain found in many SUMO E3 ligases~\citep{Hochstrasser2001}.
The SP-RING domain of Ema2 is atypical: the first cysteine of the zinc ion-binding residues found in typical counterparts is replaced by histidine, and some stabilizer residues are likely absent~\citep{Duan2009, Yunus2009}.
Despite this, GST pull-down assays demonstrated that recombinant Ema2 directly interacts with the \textit{Tetrahymena} SUMO E2 enzyme Ubc9, confirming that Ema2 is a bona fide SUMO E3 ligase.

To examine the role of Ema2 in SUMOylation in vivo, HA-tagged Smt3 (the sole SUMO gene in \textit{Tetrahymena}~\citep{Nasir2015}) was expressed in an EMA2 KO strain.
HA-Smt3 functionally replaced endogenous Smt3, as demonstrated by genetic rescue experiments in which the HA-SMT3 construct (but not the corresponding empty vector) restored viability to SMT3 KO homozygous homokaryon cells.
Similarly, a His-SMT3 construct also rescued SMT3 KO lethality, while the corresponding empty vectors (pBCMB1-His and pBCS3B1-HA) failed to restore viability in two independent transformation experiments each.
The HA-Smt3-expressing EMA2 KO strain was crossed with either a wild-type strain (WT cross) or another EMA2 KO strain (KO cross).

Western blotting using an anti-HA antibody revealed that at both 4.5 and 6~hpm, high molecular weight proteins (mainly $>$200~kDa) were detected in the WT cross but were reduced to ${\sim}50\%$ in the EMA2 KO cross.
Signal intensities of anti-HA blots in entire individual lanes were quantified across three independent experiments.
Values in the WT cross were normalized to 1, and means and standard deviations are presented with $p$-values determined by the Welch two-sample $t$-test.
These results indicate that Ema2 is the major SUMO E3 ligase during mid-conjugation.
The residual Ema2-independent SUMOylation is likely mediated by other SP-RING-containing proteins (TTHERM\_00227730, TTHERM\_00442270, TTHERM\_00348490) and/or E3-independent SUMOylation~\citep{Sampson2001}.
The requirement of SUMOylation for DNA elimination in ciliates was previously demonstrated by knockdown of UBA2 and SUMO in \textit{Paramecium}~\citep{Matsuda2006}.

\subsection{Ema2 Is Required for SUMOylation of Spt6}

To identify Ema2-dependent SUMOylation targets, His-tagged Smt3 (which also replaced Smt3 function) was expressed in an EMA2 KO strain and crossed with a wild-type (WT cross) or another EMA2 KO strain (KO cross).
SUMOylated proteins at 6~hpm were concentrated using Ni-NTA beads and identified by mass spectrometry.
Proteins from wild-type crosses without His-Smt3 expression were used to exclude non-specific binding (log2 label-free quantification [LFQ] score $>$25 or those with $>$6 consecutive histidine residues).

Spt6, the most abundantly detected protein in the WT cross, was detected at a much lower level in the EMA2 KO cross, whereas most other proteins were detected similarly between conditions.
This result identifies Spt6 as the major Ema2-dependent SUMOylation target.

To confirm this finding, HA-tagged Spt6 (Spt6-HA) was expressed from the endogenous SPT6 locus in an EMA2 KO background.
Western blotting at 4.5 and 6~hpm revealed slower-migrating Spt6-HA species in the WT cross that were barely detectable in the EMA2 KO cross.
Immunoprecipitation of Spt6-HA followed by western blotting with anti-Smt3 antibody confirmed that the slower-migrating species represent SUMOylated Spt6, which was greatly reduced in the absence of Ema2.
Consistent with the timing of Ema2 expression, these slower-migrating Spt6-HA species were absent from vegetative and starved cells (0~hpm) and were also detected at 8~hpm when the new MAC had formed.

\subsection{Ema2 Is Required for Accumulation of lncRNA in the Parental MAC}

Spt6 is a conserved regulator of multiple transcription steps across eukaryotes~\citep{Garber2023}.
Because TDSD was proposed to depend on the base-pairing interaction between scnRNAs and nascent pMAC-lncRNAs~\citep{Aronica2008, Noto2018}, we hypothesized that Ema2-dependent Spt6 SUMOylation promotes pMAC-lncRNA transcription.

RT-PCR assays amplifying transcripts spanning IES--MDS borders (which distinguish MAC-lncRNAs from MIC-lncRNAs by size) were performed at three genomic loci (M, L8, and R2).
For all three loci, pMAC-lncRNAs were detected in wild-type cells but not in EMA2 KO cells at 6~hpm, although the control RPL21 mRNA was detected in both conditions.
Genomic PCR confirmed that the corresponding MAC loci were intact in EMA2 KO cells, ruling out primer binding site loss through prior DNA elimination events.

Immunostaining using the long double-stranded RNA (dsRNA)-specific J2 antibody~\citep{Schoenborn1991} provided cytological confirmation.
In wild-type cells, long dsRNAs (reflecting bidirectional lncRNA transcription) were detected in the MIC at 3~hpm (premeiotic stage), the parental MAC at 6~hpm (mid-conjugation), and the new MAC at 9~hpm (late conjugation).
In EMA2 KO cells, dsRNAs were detected in the MIC at 3~hpm and the new MAC at 9~hpm but were undetectable in the parental MAC at 6~hpm, indicating that Ema2 is specifically required for pMAC-lncRNA accumulation.

In contrast, EMA1 KO cells showed normal dsRNA accumulation in the parental MAC, consistent with the established role of Ema1 in promoting lncRNA--scnRNA interaction rather than lncRNA transcription~\citep{Aronica2008}.
Notably, the less severe TDSD defect observed in EMA2 KO compared to EMA1 KO cells indicates that some Ema1-dependent TDSD may be initiated by single-stranded lncRNAs or mRNAs transcribed independently of Ema2.

In parallel, H3K27me3 deposition (catalyzed by Ezl1 within PRC2) was undetectable in the parental MAC of EMA2 KO cells at 6~hpm, mirroring the phenotype of EMA1, TWI1, and EZL1 KO strains~\citep{Liu2007}.
This can be explained by loss of PRC2 chromatin recruitment through the scnRNA--lncRNA interaction (Table~\ref{tab:phenotype_summary}).

\begin{table}[htbp]
\centering
\caption{Summary of phenotypic analyses across mutant strains. Phenotypes assessed in the parental MAC at mid-conjugation (4.5--6~hpm).}
\label{tab:phenotype_summary}
\begin{tabular}{lccccc}
\toprule
\textbf{Phenotype} & \textbf{WT} & \textbf{EMA2 KO} & \textbf{EMA1 KO} & \textbf{TWI1 KO} & \textbf{EZL1 KO} \\
\midrule
pMAC-lncRNA (RT-PCR, 3 loci) & Detected & Not detected & N/T & N/T & N/T \\
dsRNA (J2 staining, pMAC) & Present & Absent & Present & N/T & N/T \\
H3K27me3 (pMAC) & Present & Absent & Absent & Absent & Absent \\
TDSD (MDS scnRNA reduction) & Complete & Partial & Blocked & N/T & N/T \\
DNA elimination & Complete & Partial & N/T & Blocked & N/T \\
Progeny viability & Viable & Lethal & N/T & Lethal & N/T \\
\bottomrule
\multicolumn{6}{l}{\footnotesize N/T, not tested in this study; pMAC, parental macronucleus.}
\end{tabular}
\end{table}

\subsection{Ema2 Facilitates Chromatin Association of Spt6 and RNA Polymerase II}

The localization of Spt6-HA was examined by immunofluorescence staining.
In both WT and EMA2 KO crosses, Spt6-HA was detected in the MAC at 4.5~hpm and in the new MAC at 9~hpm.
Similarly, the MAC localization of Rpb3, the third-largest subunit of RNAPII, was not affected by the absence of Ema2.
Therefore, at the cytological level, the absence of Ema2 does not alter the gross nuclear localization of Spt6 or RNAPII.

However, cell fractionation experiments revealed critical differences at the subnuclear level.
Conjugating cells at 4.5~hpm were first lysed to extract soluble cyto- and nucleoplasmic proteins (S1 fraction), and the insoluble fraction was then sonicated to release chromatin-bound proteins (S2 fraction)~\citep{Ali2018}.
In the WT cross, Spt6-HA was detected in both the S1 and S2 fractions, but the slower-migrating SUMOylated Spt6-HA was detected exclusively in the chromatin-bound S2 fraction, indicating that SUMOylated Spt6 is associated with chromatin.
A similar distribution was observed for Rpb3.
Importantly, in the EMA2 KO cross, both Spt6-HA and Rpb3 in the chromatin-bound S2 fraction were greatly reduced, a finding validated in two independent replicate experiments.

The chromatin association of Spt6 and RNAPII does not require RNA, as they were retained in the insoluble fraction after RNase A treatment without sonication, while Twi1 (which interacts with chromatin via nascent pMAC-lncRNA~\citep{Aronica2008}) was largely released into the soluble fraction after such treatment.
We conclude that Ema2 promotes the association of Spt6 and RNAPII with chromatin in the parental MAC at the mid-conjugation stage.

\subsection{Spt6 SUMOylation Is Dispensable for Parental MAC lncRNA Transcription}

To directly assess the importance of Spt6 SUMOylation, we generated Spt6 mutants with lysine-to-arginine substitutions in different domains.
\textit{Tetrahymena} Spt6 comprises an N-terminal domain of low complexity (DOL) and a downstream conserved region containing TEX-like and SH2 domains.
Four mutant constructs were produced: DOL-KR (all lysine residues in the DOL), N-KR, M-KR, and C-KR (surface lysine residues in the respective non-DOL regions, including those predicted to be SUMOylation targets by the JASSA algorithm~\citep{Beauclair2015}).

All mutant constructs rescued the lethality of SPT6 germline KO cells.
However, N-KR and M-KR cells showed severe meiotic progression and mating initiation defects, respectively, precluding SUMOylation analysis during conjugation.
DOL-KR and C-KR cells showed apparently normal conjugation progression.
Western blotting at 4.5~hpm revealed that Spt6-DOL-KR was SUMOylated similarly to wild-type Spt6, while SUMOylation was undetectable on Spt6-C-KR (Table~\ref{tab:spt6_mutants}).
Thus, Spt6-C-KR represents a SUMOylation-defective Spt6 mutant with at least reduced SUMOylation compared to Spt6 in the absence of Ema2.

\begin{table}[htbp]
\centering
\caption{Summary of Spt6 mutant analyses. Lysine-to-arginine substitutions were introduced into different Spt6 domains.}
\label{tab:spt6_mutants}
\begin{tabular}{lcccccc}
\toprule
\textbf{Construct} & \textbf{Genetic} & \textbf{Conjugation} & \textbf{SUMO-} & \textbf{Chromatin} & \textbf{pMAC} & \textbf{DNA elim.} \\
 & \textbf{rescue} & \textbf{progression} & \textbf{ylation} & \textbf{assoc.} & \textbf{lncRNA} & \textbf{defect} \\
\midrule
HA-SPT6-WT & Yes & Normal & Normal & Normal & Normal & None \\
HA-SPT6-DOL-KR & Yes & Normal & Normal & N/T & N/T & N/T \\
HA-SPT6-N-KR & Yes & Meiosis defect & N/T & N/T & N/T & N/T \\
HA-SPT6-M-KR & Yes & Mating defect & N/T & N/T & N/T & N/T \\
HA-SPT6-C-KR & Yes & Normal$^{a}$ & Absent & Normal$^{b}$ & Normal & Mild \\
\bottomrule
\multicolumn{7}{l}{\footnotesize $^{a}$Lower mating efficiency observed. $^{b}$Confirmed in three independent experiments.} \\
\multicolumn{7}{l}{\footnotesize N/T, not tested.}
\end{tabular}
\end{table}

Critically, cell fractionation showed that Spt6-HA and RNAPII were present in the chromatin-bound fraction of Spt6-C-KR cells at levels comparable to wild-type, consistently across three independent experiments.
Furthermore, J2 antibody staining detected long dsRNAs in the parental MAC (and the new MAC) of Spt6-C-KR cells.
These results demonstrate that, contrary to initial expectations, Spt6 SUMOylation per se is not required for pMAC-lncRNA transcription.
Nevertheless, exconjugants from Spt6-C-KR cells showed a mild DNA elimination defect, as detected by FISH, indicating that Ema2-directed Spt6 SUMOylation plays a role in DNA elimination that is distinct from promoting pMAC-lncRNA transcription.

%% ============================================================
%% 3. DISCUSSION
%% ============================================================
\section{Discussion}

In this study, we demonstrated that the conjugation-specific SUMO E3 ligase Ema2 is required for the accumulation of lncRNAs, TDSD, and heterochromatin formation in the parental MAC, and ultimately for completing programmed DNA elimination in \textit{Tetrahymena}.
Ema2 is responsible for SUMOylation of the transcriptional regulator Spt6 and promotes the interaction of Spt6 and RNAPII with chromatin.
We therefore conclude that Ema2 facilitates genome-wide lncRNA transcription in the parental MAC, which is a prerequisite for scnRNA--chromatin communication and the downstream TDSD that regulates DNA elimination.

\subsection{Ema2-Dependent Chromatin Association of Transcriptional Machinery}

The observation that Ema2 enhances the chromatin interaction of Spt6 and RNAPII appears to contradict the fact that Spt6 and RNAPII are essential for vegetative cell viability~\citep{Mochizuki2004b}, while Ema2 is expressed exclusively during conjugation.
Moreover, EMA2 KO cells did not significantly impede the progression of conjugation processes, indicating that essential mRNA transcription continues in the parental MAC even without Ema2.
Therefore, the loss of the majority of Spt6 and RNAPII from chromatin in the absence of Ema2 must be a temporal event during the mid-conjugation stage.
This suggests that Spt6 and RNAPII might be specifically engaged in pMAC-lncRNA transcription at this particular time window in wild-type cells.

It is likely that radical changes in the chromatin environment of the parental MAC necessitate Ema2-dependent SUMOylation to maintain Spt6 and RNAPII on chromatin for global pMAC-lncRNA transcription.
Alternatively, because expression of many transcriptional machineries, including Spt6 and RNAPII subunits, is highly upregulated during conjugation~\citep{Miao2009, Xiong2012, Mochizuki2004b}, pMAC-lncRNA transcription may require a greater amount of transcriptional machinery than other transcription types, and Ema2-dependent SUMOylation may be required for committing these excess machineries to chromatin.

\subsection{Spt6 SUMOylation Is Dispensable for lncRNA Transcription}

The finding that the SUMOylation-defective Spt6-C-KR mutant still supported pMAC-lncRNA transcription and chromatin association suggests that Ema2 may facilitate lncRNA transcription through SUMOylation of protein(s) other than Spt6.
Although mass spectrometry did not detect other Ema2-dependent SUMOylation events, such events may be masked by Ema2-independent SUMOylation of the same proteins or by SUMOylation-triggered protein degradation through SUMO-targeted ubiquitination~\citep{Praefcke2012, Staudinger2017}.
Alternatively, SUMOylation of Spt6 may promote pMAC-lncRNA transcription by competing with another modification at the same lysine residues, and the Spt6-C-KR mutant might actually mimic the SUMOylated state by preventing such modification.
The mild DNA elimination defect observed in Spt6-C-KR exconjugants suggests that Spt6 SUMOylation plays a distinct role in DNA elimination beyond promoting pMAC-lncRNA transcription.

\subsection{SUMOylation in Small RNA-Directed Chromatin Regulation Across Species}

SUMOylation has been implicated in small RNA-directed chromatin regulation in multiple organisms.
In fission yeast, SUMOylation of the heterochromatin components Swi6, Chp2, and Clr4 is important for RNAi-directed heterochromatin silencing~\citep{Shin2005}.
In \textit{C. elegans}, SUMOylation of HDAC1 promotes Argonaute-directed transcriptional silencing~\citep{Kim2021}.
In \textit{Drosophila} piRNA-directed transposon silencing, the SUMO E3 ligase Su(var)2-10 recruits the SetDB1 histone methyltransferase complex for H3K9me3 deposition~\citep{Ninova2020a, Ninova2020b}, and Panoramix SUMOylation recruits the heterochromatin effector Sov~\citep{Andreev2022}.
All these SUMOylation events appear to require pre-existing lncRNA transcription.
In contrast, the present study shows that Ema2-dependent SUMOylation during DNA elimination in \textit{Tetrahymena} operates upstream of lncRNA transcription.

Furthermore, the loss of the Piwi protein Twi1 blocks H3K27me3 accumulation in the parental MAC~\citep{Liu2007, Xu2021} without affecting lncRNA accumulation~\citep{Woo2016}, indicating that heterochromatin formation is not a prerequisite for lncRNA transcription in \textit{Tetrahymena}.
This contrasts with the heterochromatin-dependent lncRNA transcription of piRNA clusters in \textit{Drosophila}~\citep{Andersen2017}.
\textit{Tetrahymena} therefore likely employs a unique evolutionary solution for transcribing lncRNAs to promote TDSD and small RNA-directed heterochromatin formation.
Further investigation of Ema2-dependent SUMOylation in lncRNA transcription will clarify an undescribed regulatory layer in small RNA-directed chromatin regulation.

%% ============================================================
%% 4. MATERIALS AND METHODS
%% ============================================================
\section{Materials and Methods}

\subsection{Strains and Cell Culture}

The \textit{Tetrahymena thermophila} wild-type strains B2086, CU427, and CU428 and the MIC-defective ``star'' strains B*VI and B*VII were obtained from the Tetrahymena Stock Center.
The production of EMA2 KO strains was described previously~\citep{Shehzada2022}.
Cells were grown in SPP medium~\citep{Gorovsky1975} containing 2\% proteose peptone at 30$^{\circ}$C.
For conjugation, growing cells (${\sim}5$--$7 \times 10^{5}$/mL) of two different mating types were washed, pre-starved (${\sim}12$--24~hr), and mixed in 10~mM Tris (pH~7.5) at 30$^{\circ}$C.

\subsection{Establishment of EMA2-HA Strains}

The C-terminal coding sequence of EMA2 and the 3$'$ flanking sequence were amplified by PCR from B2086 genomic DNA, combined by overlapping PCR, and cloned into pBlueScriptSK(+).
The HA-neo3 cassette~\citep{Kataoka2010} was inserted to generate pEMA2-HA-neo3.
This construct was introduced into the EMA2 MAC locus by biolistic transformation~\citep{Bruns2000}, followed by selection with 100~$\mu$g/mL paromomycin and 1~$\mu$g/mL CdCl$_{2}$.
Endogenous EMA2 MAC copies were replaced through phenotypic assortment~\citep{Hamilton2000}.

\subsection{DNA Elimination Assay by FISH}

DNA elimination was analyzed by FISH as described~\citep{Loidl2004}.
Plasmid DNAs pMBR 4C1, pMBR 2, and Tlr IntB, containing different parts of the Tlr1 element~\citep{Wuitschick2002}, were used to produce Cy5-labeled FISH probes by nick translation.
The IES retention index was calculated as:
\begin{equation}
\text{IES retention index} = 0.35 \times \frac{(\text{mean Cy3}_A - \text{mean Cy3}_C) / (\text{mean Cy3}_I - \text{mean Cy3}_C)}{(\text{mean DAPI}_A - \text{mean DAPI}_C) / (\text{mean DAPI}_I - \text{mean DAPI}_C)}
\end{equation}
where $A$, $I$, and $C$ denote the new MAC, MIC, and cytoplasm areas, respectively.
The correction factor of 0.35 was obtained experimentally to set the average IES retention index of TWI1 KO cells to 1, based on the observation that DNA elimination is completely blocked in TWI1 KO.

\subsection{Small RNA Analyses}

For northern hybridization, 10~$\mu$g of total RNA from conjugating wild-type and EMA2 KO cells (isolated using TRIzol reagent) was separated in a 15\% acrylamide-urea gel and analyzed by northern blot using the radio-labeled Mi-9 probe as described~\citep{Aronica2008}.
High-throughput sequencing of small RNAs was performed as described~\citep{Mutazono2019}.
The 2020 MAC genome assembly~\citep{Sheng2020} was fragmented into 10~kb pieces, and each fragment containing $>$3~kb of mappable sequence was used as an MDS tile (total 10,235 tiles).
Each Type-A IES~\citep{Noto2015} served as an IES tile (total 4,691 tiles).
Normalized numbers (RPKM [reads per kilobase of unique sequences per million]) of 26--32-nt small RNAs uniquely matching tiles were counted.

\subsection{GST Pull-Down Assay}

\textit{E. coli} codon-optimized EMA2 and UBC9 genes were cloned into pGEX-4T-1 and pET28a(+), respectively.
His-Ubc9 was expressed in BL21 \textit{E. coli} (${\sim}$16~hours at 25$^{\circ}$C with 0.05~mM IPTG), purified using Ni-NTA beads (lysis in 50~mM HEPES pH~7.5, 800~mM KCl, 10\% glycerol, 0.2~mM PMSF, 1$\times$ cOmplete EDTA-free; washes with 20~mM imidazole; elution with 250~mM imidazole), dialyzed, and concentrated through a Microcon 10~kDa filter.
GST-Ema2 or GST alone was expressed and immobilized on Glutathione Sepharose 4B (90~min at 4$^{\circ}$C).
Beads were incubated with 10~$\mu$g His-Ubc9 in pull-down buffer (20~mM Tris-HCl pH~7.5, 100~mM NaCl, 1\% NP-40, 5~mM MgCl$_{2}$, 1~mM DTT, 0.2~mM PMSF, 1$\times$ cOmplete EDTA-free) for 90~min at 4$^{\circ}$C, washed 4 times, and eluted with 2$\times$ SDS sample buffer.

\subsection{Purification of SUMOylated Proteins}

For His-Smt3-conjugated protein purification, cells were lysed in Guanidine Lysis Buffer (6~M guanidine-HCl, 93.2~mM Na$_{2}$HPO$_{4}$, 6.8~mM NaH$_{2}$PO$_{4}$, 10~mM Tris-HCl pH~8.0) by sonication ($2 \times 25$ pulses, 20\% power, 60\% pulse).
The cleared lysate was supplemented with 50~mM imidazole and 5~mM $\beta$-mercaptoethanol and incubated with Ni-NTA agarose beads.
Bound proteins were eluted with Elution buffer (6~M urea, 58~mM Na$_{2}$HPO$_{4}$, 42~mM NaH$_{2}$PO$_{4}$, 10~mM Tris-HCl pH~8.0, 500~mM imidazole) and analyzed by western blot and mass spectrometry.

For HA-Smt3-conjugated protein purification, total proteins were precipitated with 6\% TCA on ice for 10~min, pelleted at 13,000~rpm for 5~min at 4$^{\circ}$C, dissolved in 1$\times$ SDS sample buffer, neutralized with 2~M Tris, and incubated at 95$^{\circ}$C for 5~min.
The dissolved proteins (600~$\mu$L) were then suspended in 7.4~mL of binding buffer (20~mM Tris pH~7.5, 100~mM NaCl, 2~mM MgCl$_{2}$, 2~mM CaCl$_{2}$) and incubated with EZview Red anti-HA agarose beads (Sigma) overnight at 4$^{\circ}$C.
Beads were washed 5 times with wash buffer (20~mM Tris pH~7.5, 100~mM NaCl, 2~mM MgCl$_{2}$, 2~mM CaCl$_{2}$, 0.1\% Tween~20) and eluted with 1$\times$ SDS buffer at 95$^{\circ}$C for 5~min.

\subsection{Functionality Test of HA-SMT3 and His-SMT3}

Two SMT3 homozygous heterokaryon KO strains (B1-1 and B5-17) were mated, and pBCMB1-His-SMT3 (or the empty pBCMB1-His vector) or pBCS3B1-HA-SMT3 (or the empty pBCS3B1-HA vector) constructs were introduced into the new MAC at 8~hpm by particle gun.
Cells were incubated in 10~mM Tris pH~7.5 at 30$^{\circ}$C until 24~hpm and fed by adding an equal volume of 2$\times$ SPP.
After 3~hr incubation at 30$^{\circ}$C, neo expression (and His-Smt3 expression for pBCMB1-His-SMT3) was induced by adding 1~$\mu$g/mL CdCl$_{2}$, and cells were incubated for an additional 1~hr.
Then, 100~$\mu$g/mL paromomycin was added and cells were aliquoted to 96-well plates (150~$\mu$L/well, 7 plates).
Plates were incubated at 30$^{\circ}$C for 4--5 days, and paromomycin-resistant cells were confirmed by genomic PCR.
For the empty vector controls, two independent transformation experiments were performed to confirm their inability to restore cell viability of KO cells.

\subsection{Cell Fractionation}

Cell fractionation was performed as described~\citep{Ali2018}.
A total of $1.4 \times 10^{6}$ cells were incubated with 1~mL of Lysis buffer (50~mM Tris pH~8.0, 100~mM NaCl, 5~mM MgCl$_{2}$, 1~mM EDTA, 0.5\% Triton X-100, 1$\times$ Complete Protease Inhibitor) on ice for 30~min and centrifuged at 15,000$g$ for 10~min at 4$^{\circ}$C.
Soluble proteins (S1 fraction; 500~$\mu$L) were precipitated with 10\% TCA and dissolved in 70~$\mu$L of 1$\times$ SDS sample buffer.
The precipitate was resuspended in 2~mL of fresh Lysis buffer, sonicated (5 rounds of 20\% power, 50\% pulse, 5 pulses each), and centrifuged at 20,000$g$ for 10~min at 4$^{\circ}$C.
The supernatant (1~mL; chromatin-bound S2 fraction) was precipitated with 10\% TCA and dissolved in 70~$\mu$L of 1$\times$ SDS sample buffer.

For RNase treatment experiments, $1.4 \times 10^{6}$ cells were incubated with 1~mL of Lysis buffer containing 20~$\mu$g/mL RNase A on ice for 30~min.
The soluble (S1), sonication-solubilized (S2), and insoluble (I) fractions were obtained and analyzed by western blotting.

\subsection{RT-PCR Analysis of lncRNA Transcripts}

Total RNA from mating wild-type and EMA2 KO cells was isolated at 6~hpm using TRIzol reagent.
lncRNA transcripts from MAC loci containing the M-, L8-, and R2 IES boundaries were amplified by nested PCR (95$^{\circ}$C/20~s, 50$^{\circ}$C/30~s, 68$^{\circ}$C/1~min; 30 cycles each) and analyzed on 1\% agarose gels.
RPL21 mRNA served as a positive control.

\subsection{Establishment of SPT6 Germline KO and Rescue Strains}

The knockout construct for SPT6 was generated by amplifying and concatenating the 5$'$ and 3$'$ flanking regions by overlap PCR.
For genetic rescue, constructs expressing HA-tagged wild-type Spt6 (HA-SPT6-WT) or K-to-R mutants (DOL-KR, N-KR, M-KR, C-KR) were introduced into SPT6 KO homozygous homokaryon cells.
Transformed cells were selected with 300~$\mu$g/mL puromycin and 1~$\mu$g/mL CdCl$_{2}$ and phenotypically assorted as described~\citep{Hamilton2000} until cells grew in 1.2~mg/mL puromycin.

\subsection{Immunofluorescence Staining and Western Blotting}

Immunofluorescence staining of long dsRNA was performed as described~\citep{Woo2016}.
Other immunofluorescence experiments were performed as described~\citep{Loidl2004}.
All primary antibodies and respective secondary antibodies (Alexa-488-coupled anti-rabbit or mouse IgG, Invitrogen) were diluted 1:1,000.
For western blotting, all primary and respective secondary antibodies (HRP-coupled anti-rabbit or mouse IgG, Jackson ImmunoResearch Lab) were diluted 1:10,000.
The following antibodies were used: anti-H3K27me3 (Millipore, 07-449), anti-HA HA11 (Covance, clone 6B12), anti-GST (BD Biosciences, 554805), anti-His (Proteintech, 66005-1-Ig), anti-long dsRNA J2 (Jena Bioscience), anti-$\alpha$-tubulin 12G10 (Developmental Studies Hybridoma Bank), anti-Rpb3~\citep{Kataoka2017}, and anti-Smt3~\citep{Nasir2015}.

\subsection{Progeny Viability Test}

Mating pairs were isolated into drops of 1$\times$ SPP at 6~hpm and incubated for ${\sim}3$ days.
Completion of conjugation was assessed by 6-methylpurine resistance and paromomycin sensitivity, as described~\citep{Mochizuki2002}.

%% ============================================================
%% BIBLIOGRAPHY
%% ============================================================
\bibliography{references}

\end{document}
